% =========================================================================
\section{Introducción}\label{sec:introduccion}
% =========================================================================
En esta sesión montamos con un $\mu$A741, alimentado a $\pm10\,\mathrm{V}$, los circuitos diseñados en el PreLab:
un amplificador inversor de ganancia $-10$, uno no inversor de ganancia $+11$ y, después, un integrador y un
derivador inversores. Con el osciloscopio medimos las formas de onda de entrada y salida y las
características de transferencia en modo X--Y, y las comparamos con lo que predice el modelo ideal del
amplificador operacional. Este informe recoge esas medidas y responde, en orden, a las cuestiones del guion.

\begin{center}\small
\renewcommand{\arraystretch}{1.3}
\begin{tabular}{@{}lll@{}}
\toprule
Integrante & Apellidos, nombre & DNI/Pasaporte\\
\midrule
1 & \pend{\phantom{Apellido Apellido, Nombre}} & \pend{\phantom{00000000X}}\\
2 & \pend{\phantom{Apellido Apellido, Nombre}} & \pend{\phantom{00000000X}}\\
\bottomrule
\end{tabular}
\end{center}

\textbf{Componentes.} Como pide el guion, medimos con el multímetro el valor real de cada componente pasivo;
los cálculos de este informe usan los valores medidos.
\begin{center}\small
\renewcommand{\arraystretch}{1.3}
\begin{tabular}{@{}lcc@{}}
\toprule
Componente & Valor nominal & Valor medido\\
\midrule
$R_1$ (circuitos 1 y 2) & $1\,\mathrm{k}\Omega$ & \dato\\
$R_2$ (circuitos 1 y 2) & $10\,\mathrm{k}\Omega$ & \dato\\
$R_L$ & $10\,\mathrm{k}\Omega$ & \dato\\
$R$ (circuito 3) & $1\,\mathrm{k}\Omega$ & \dato\\
$R_F$ (circuito 3, set 1) & $100\,\mathrm{k}\Omega$ & \dato\\
$C$ (circuito 3) & $1{,}59\,\mathrm{nF}$ & \dato\\
Tensiones de alimentación $\pm\Vcc$ & $\pm10\,\mathrm{V}$ & \dato\\
\bottomrule
\end{tabular}
\end{center}
